\documentclass[aspectratio=169,10pt]{beamer}
\usepackage{teaching_slides}

\title{Production Functions, Productivity, and Markups}
\author{C.Conlon}
\institute{Grad IO}
\date{Fall 2026}

\begin{document}

\frame{\titlepage}

\begin{frame}{Overview}
\begin{itemize}
\item Why we estimate production functions: productivity facts, and two uses of $f(\cdot)$
\item The identification problem: simultaneity, selection, and prices (TFPR vs TFPQ; Foster--Haltiwanger--Syverson)
\item One equation, three families of estimators
\begin{itemize}
\item dynamic panel (Blundell--Bond)
\item proxy / control function: Olley--Pakes, Levinsohn--Petrin, Ackerberg--Caves--Frazer
\item first-order conditions for gross output: Gandhi--Navarro--Rivers; plus Leontief in materials
\end{itemize}
\item Endogenous productivity: exporting and R\&D belong in the law of motion
\item Markups from production data (De Loecker--Warzynski; multi-product firms), and where it breaks
\begin{itemize}
\item revenue data; invertibility; market power; heterogeneous technology; input market power
\end{itemize}
\item Aggregation (which costs, which average) and misallocation
\end{itemize}
\vfill
{\footnotesize Notation: producer $j$, period $t$. Lowercase letters are logs of uppercase levels, \alert{except prices}: output price $p_{jt}$ and input prices $w^X_{jt}$ are levels, as in the rest of the course.}
\end{frame}


\section{Why Production Functions?}

\begin{frame}{Two Uses of a Production Function}
\begin{align*}
y_{jt} = f(\mathbf{x}_{jt}) + \omega_{jt} + \varepsilon_{jt}, \qquad \mathbf{x}_{jt} = (k_{jt}, l_{jt}, m_{jt})
\end{align*}
\begin{enumerate}
\item \alert{Productivity} $\omega_{jt}$: what is left once inputs are accounted for.
\begin{itemize}
\item dispersion and persistence; reallocation and entry/exit
\item effects of deregulation, trade liberalization, exporting, R\&D, management
\end{itemize}
\item \alert{Output elasticities} $\theta^X_{jt} = \partial f / \partial x_{jt}$.
\begin{itemize}
\item returns to scale, factor-biased technology
\item \alert{markups}: $\mu_{jt} = \theta^X_{jt}/\alpha^X_{jt}$ with $\alpha^X$ the revenue share of a flexible input (De Loecker--Warzynski 2012)
\end{itemize}
\end{enumerate}
Questions this lecture can answer:
\begin{itemize}
\item How much of aggregate productivity growth comes from reallocation and turnover?
\item What can production data alone say about market power, and about its trend?
\end{itemize}
\end{frame}


\begin{frame}{Facts About Producer-Level Productivity}
\begin{itemize}
\item \alert{Dispersion}: the 90--10 TFP ratio is about 1.9:1 within US 4-digit manufacturing (log gap 0.65, Syverson 2004), about 2:1 within 6-digit industries in the US and Canada, and above 5:1 (revenue-based) in China and India (Hsieh--Klenow 2009).
\item It is not missing inputs: adding R\&D stocks or detailed worker quality barely moves it (Fox--Smeets 2011: 90--10 ratio 3.74 $\to$ 3.36).
\item \alert{Persistence}: annual AR(1) coefficients of 0.6--0.8; 0.3--0.4 with producer fixed effects.
\item \alert{Correlates}: more productive producers are larger, more profitable, grow faster, are more likely to survive, pay higher wages, and (for homogeneous goods) charge lower prices.
\end{itemize}
\vfill
\begin{quote}
``I found the spectacle of economic models yielding large residuals rather uncomfortable, even when the issue was fudged by renaming them technical change and claiming credit for their `measurement.'\,'' --- Zvi Griliches
\end{quote}
{\footnotesize Sources: Syverson (2011, JEL); De Loecker and Syverson (2021, Handbook of IO).}
\end{frame}


\begin{frame}{What Is Productivity?}
Hicks-neutral productivity $\Omega$: \quad $Y = \Omega\, F(K, L, M)$.
\begin{itemize}
\item Three equivalent views (De Loecker--Syverson 2021):
\begin{itemize}
\item a shifter of the isoquants;
\item an output-to-input ratio, $\Omega = Y/F(\cdot)$: total factor productivity (TFP);
\item a shifter of the cost curve: higher $\Omega$, lower cost at every $Y$.
\end{itemize}
\item Labor productivity $Y/L$ moves with $K/L$ even when technology does not, so we usually want TFP. TFP needs a functional form and a measure of capital.
\item Data usually record revenue, not quantity. In logs:
\begin{align*}
\text{TFPR}_{jt} = \text{TFPQ}_{jt} + \ln p_{jt}.
\end{align*}
\item TFPR is \alert{not} the residual of a production function estimated with revenue on the left. We derive the gap below, after Klette--Griliches.
\end{itemize}
\end{frame}


\begin{frame}{What Moves Productivity? (Syverson 2011)}
\begin{columns}[T]
\begin{column}{0.5\textwidth}
\textbf{Inside the producer}
\begin{itemize}
\item Management practices: the interquartile range of management scores is worth 3.2--7.5\% of TFP (Bloom--Van Reenen 2007).
\item Input quality: worker skill, capital vintage.
\item IT and R\&D (Doraszelski--Jaumandreu 2013: R\&D also makes productivity riskier).
\item Learning by doing: Lockheed L-1011 hours per plane halved by the 30th plane (Benkard 2000).
\item Product innovation; firm structure.
\end{itemize}
\end{column}
\begin{column}{0.5\textwidth}
\textbf{Outside the producer}
\begin{itemize}
\item Spillovers (Liberty ships).
\item Competition: US iron ore labor productivity doubled 1980--85 under import pressure (Schmitz 2005).
\item Deregulation: telecom equipment (Olley--Pakes 1996).
\item Trade; flexible input markets.
\item Misallocation across producers (Hsieh--Klenow 2009; later).
\end{itemize}
\end{column}
\end{columns}
\vfill
All of these need a credible estimate of $\omega_{jt}$, which needs a credible estimate of $f$. Firm choices among them (R\&D, exporting) belong in the law of motion for $\omega$ (later).
\end{frame}


\section{The Identification Problem}

\begin{frame}{Setup and Timing}
\begin{align*}
y_{jt} = f(k_{jt}, l_{jt}, m_{jt}) + \omega_{jt} + \varepsilon_{jt}, \qquad
\omega_{jt} = g(\omega_{jt-1}) + \xi_{jt}, \qquad \E[\xi_{jt} \mid \mathcal{I}_{jt-1}] = 0.
\end{align*}
\begin{itemize}
\item $\omega_{jt}$ is known to the firm when it chooses period-$t$ inputs ($\omega_{jt} \in \mathcal{I}_{jt}$); $\varepsilon_{jt}$ is not (an ex post shock or measurement error).
\item First-order Markov productivity. In this lecture \alert{$\xi_{jt}$ is the productivity innovation}, not BLP's demand unobservable; demand shifters will be written $\delta_{jt}$. With an AR(1), $g(\omega) = \rho\,\omega$ ($\rho$ = persistence, not the nesting parameter).
\item Input timing (the GNR taxonomy):
\begin{itemize}
\item \alert{predetermined}: $k_{jt} \in \mathcal{I}_{jt-1}$, chosen through investment $i_{jt-1}$;
\item \alert{flexible}: chosen at $t$ after seeing $\omega_{jt}$, no dynamic implications (materials; often labor);
\item \alert{dynamic}: lagged choices affect current ones (adjustment costs).
\end{itemize}
\item Data: many producers, few periods ($N \to \infty$, $T$ fixed).
\end{itemize}
\end{frame}


\begin{frame}{Simultaneity (Marschak and Andrews 1944)}
\small
Regress output on inputs and the slope is $\theta^X$ plus a \alert{transmission bias}, $\partial \E[\omega_{jt} \mid \mathbf{x}_{jt}]/\partial x$. To see its size, take Cobb--Douglas value added with $k$ predetermined, $l$ flexible, price taking, and a firm-level price ratio $d_{jt} \equiv \ln(w^L_{jt}/p_{jt})$ independent of $\omega$. The labor FOC gives
\begin{align*}
l_{jt} = \frac{\ln\beta_l - d_{jt} + \beta_k k_{jt} + \omega_{jt}}{1 - \beta_l}
\quad\Longrightarrow\quad
\operatorname{plim} \hat\beta_l^{\,OLS} = \beta_l + (1 - \beta_l)\, \frac{\sigma^2_\omega}{\sigma^2_\omega + \sigma^2_d},
\end{align*}
with variances net of $k$.
\begin{itemize}
\item With no price dispersion ($\sigma^2_d = 0$), $\omega$ is an exact function of $(l, k)$ and $y_{jt} = l_{jt} + \text{const} + \varepsilon_{jt}$: OLS puts 1 on labor and 0 on capital, whatever the truth.
\item The more an input responds to $\omega$ rather than to prices, the closer its coefficient gets to 1. In data, OLS overstates the materials elasticity by 34\% (Chile and Colombia, GNR 2020) and labor in Olley--Pakes's data. Biases on less flexible inputs can go either way.
\item $\sigma^2_d$ is the variation an input-price IV would use, but firm-level prices are rarely observed and often reflect quality or an upward-sloping supply curve (Griliches--Mairesse 1998).
\item $\sigma^2_d = 0$ comes back twice: in LP's first stage (ACF) and for gross output (GNR).
\end{itemize}
\end{frame}


\begin{frame}{Selection (Olley and Pakes 1996)}
Firms exit when productivity falls below a threshold that depends on their state:
\begin{align*}
\chi_{jt} = \mathbf{1}\{\omega_{jt} \ge \underline{\omega}_t(k_{jt}, a_{jt})\}.
\end{align*}
\begin{itemize}
\item Larger capital stocks mean higher continuation values, so $\underline{\omega}_t$ is \alert{decreasing} in $k$: big plants survive worse draws.
\item Among survivors, $\E[\omega_{jt} \mid k_{jt}, \chi_{jt} = 1]$ is decreasing in $k$, which biases $\hat\beta_k$ \alert{down}.
\item OP's telecom data: a balanced panel throws away about 65\% of observations and gives $\hat\beta_k = 0.173$; the full sample with their correction gives $0.355$.
\item Same logic whenever selection is on productivity: exporters (Melitz 2003), entrants, survivors of a trade shock.
\end{itemize}
\end{frame}


\begin{frame}{Measuring Output: Revenue Is Not Quantity}
Most producer data record sales and expenditures. With an industry deflator $\bar p_t$:
\begin{align*}
r_{jt} - \ln \bar p_t = f(\mathbf{x}_{jt}) + \omega_{jt} + \varepsilon_{jt} + \underbrace{(\ln p_{jt} - \ln \bar p_t)}_{\text{omitted price}}.
\end{align*}
\begin{itemize}
\item If firms set prices, $\ln p_{jt}$ falls as output rises, so it is correlated with every input. Klette and Griliches (1996): the estimated elasticities are scaled down by $1/\mu$ (next slide).
\item Input expenditures have the same problem when input prices vary across firms.
\item Even with quantities: multi-product producers, quality differences, units that are not comparable across firms.
\item We come back to this when we estimate markups, where it matters most.
\end{itemize}
\end{frame}


\begin{frame}{TFPR Is Not the Revenue Residual}
\small
Klette and Griliches (1996): $y_{jt} = f(\mathbf{x}_{jt}) + \omega_{jt}$ (drop $\varepsilon$) and CES demand $y_{jt} - \bar y_t = -\eta\,(\ln p_{jt} - \ln \bar p_t) + \delta_{jt}$, with demand shifter $\delta_{jt}$ and $\eta > 1$, so $\mu = \eta/(\eta - 1)$. Substitute the price out of deflated revenue $\tilde r_{jt} \equiv r_{jt} - \ln \bar p_t = y_{jt} + \ln p_{jt} - \ln \bar p_t$:
\begin{align*}
\tilde r_{jt} &= \underbrace{\tfrac{\eta-1}{\eta}\, f(\mathbf{x}_{jt})}_{\text{elasticities } \theta^X/\mu} + \tfrac{1}{\eta}\, \bar y_t + \underbrace{\tfrac{\eta-1}{\eta}\, \omega_{jt} + \tfrac{1}{\eta}\, \delta_{jt}}_{\omega^R_{jt}:\ \text{revenue residual}}, \\
\text{TFPR}_{jt} &\equiv \tilde r_{jt} - f(\mathbf{x}_{jt}) = \omega_{jt} + \ln p_{jt} - \ln \bar p_t = \omega^R_{jt} - \tfrac{1}{\eta}\big(f(\mathbf{x}_{jt}) - \bar y_t\big).
\end{align*}
\begin{itemize}
\item The revenue regression recovers $\theta^X/\mu$, not $\theta^X$: the Klette--Griliches bias.
\item The gap grows with the firm's inputs at rate $1/\eta$, a demand parameter: big firms look better by $\omega^R$.
\item Starkest case (Hsieh--Klenow): CRS, flexible inputs, common input prices $\mathbf{w}_t$. Then $mc = c(\mathbf{w}_t)/\Omega$ and $p = \mu\, mc$, so $\text{TFPR}_{jt} = \ln \mu + \ln c(\mathbf{w}_t) - \ln \bar p_t$: \alert{the same for every firm}. With $\eta = 4$, 10\% more TFPQ raises $\omega^R$ by 7.5\% and a 10\% demand shock raises it by 2.5\%; neither moves TFPR.
\end{itemize}
{\footnotesize De Loecker--Syverson (2021), fns.\ 6 and 42. De Loecker (2011) estimates $\eta$ inside the production function (later).}
\end{frame}


\begin{frame}{Foster, Haltiwanger and Syverson (2008)}
\small
Eleven physically homogeneous products (ready-mixed concrete, block ice, sugar, plywood, \ldots), US Census of Manufactures 1977--97, 17,669 plant-years. Prices and quantities are both observed.

Two benchmarks bracket the data. Price taking: $\operatorname{corr}(p, \text{TFPQ}) = 0$ and TFPR $=$ TFPQ. Hsieh--Klenow (previous slide): $\operatorname{corr}(p, \text{TFPQ}) = -1$ and TFPR does not vary.
\begin{columns}[T]
\begin{column}{0.5\textwidth}
\begin{center}
\scriptsize
\renewcommand{\arraystretch}{1.1}
\begin{tabular}{lccc}
\toprule
Correlations & Price & TFPR & TFPQ \\
\midrule
Price & 1.00 & & \\
TFPR & 0.16 & 1.00 & \\
TFPQ & $-0.54$ & 0.75 & 1.00 \\
\midrule
s.d. & 0.18 & 0.22 & 0.26 \\
1-yr persistence & 0.82 & 0.77 & 0.79 \\
\bottomrule
\end{tabular}

{\tiny FHS Tables 1 and 3; product-year effects removed. Demand shocks: persistence 0.91.}
\end{center}
\end{column}
\begin{column}{0.48\textwidth}
\begin{itemize}
\item The data sit in between: more physically productive plants charge \alert{lower} prices, so TFPR understates TFPQ dispersion.
\item Entrants have higher TFPQ than incumbents but charge lower prices: revenue measures understate the entrant advantage.
\item Exit depends on TFPQ \alert{and} on demand. Selection is on profitability, not just productivity.
\item Demand shocks are more persistent than productivity.
\end{itemize}
\end{column}
\end{columns}
\end{frame}


\begin{frame}{Classic Fixes and What They Assume}
\small
\renewcommand{\arraystretch}{1.25}
\begin{tabularx}{\textwidth}{@{}l X X@{}}
\toprule
Fix & Assumes & Problem in practice \\
\midrule
Fixed effects & $\omega_{jt} = \omega_j$ & tiny $\hat\beta_k$ and low returns to scale; differencing amplifies measurement error (OP within: $\hat\beta_k = 0.067$) \\
Input-price IV & firm-level $w^X$ observed and exogenous & rarely observed; quality and location confounds \\
Cost shares / index numbers & static FOCs for \emph{all} inputs, CRS, no ex post shock & capital and labor have adjustment costs; with revenue shares also needs $\mu = 1$ \\
\bottomrule
\end{tabularx}
\vspace{1em}
\begin{itemize}
\item The modern estimators relax some of these and add others: a timing assumption, a Markov process for $\omega$, and either linearity or a scalar unobservable.
\item ACF's advice: be explicit about which assumptions you are making, and whether they fit the industry.
\end{itemize}
\end{frame}


\section{One Equation, Three Families}

\begin{frame}{One Equation Behind Every Estimator}
Substitute the Markov process into the production function:
\begin{align*}
y_{jt} = f(\mathbf{x}_{jt}) + g(\omega_{jt-1}) + \xi_{jt} + \varepsilon_{jt}, \qquad \E[\xi_{jt} + \varepsilon_{jt} \mid \mathbf{z}_{jt}] = 0,
\end{align*}
with instruments $\mathbf{z}_{jt} \subseteq \mathcal{I}_{jt-1}$, e.g.\ $(k_{jt}, l_{jt-1}, m_{jt-1})$. The estimators differ in how they \alert{replace $\omega_{jt-1}$}:
\vspace{0.5em}

\small
\renewcommand{\arraystretch}{1.3}
\begin{tabularx}{\textwidth}{@{}l l X@{}}
\toprule
Family & $\omega_{jt-1}$ replaced by & Key extra assumption \\
\midrule
Dynamic panel & $y_{jt-1} - f(\mathbf{x}_{jt-1}) - \varepsilon_{jt-1}$ & $g$ linear \\
Proxy / control function & $h_{t-1}(\mathbf{o}_{jt-1})$, an inverted input demand & scalar unobservable, monotone input demand \\
FOC / share form & $y_{jt-1} - f(\mathbf{x}_{jt-1}) + \ln\big(\alpha^X_{jt-1}/\theta^X_{jt-1}\big)$ & static FOC for $X$; price taking (or its cost-share version) \\
\bottomrule
\end{tabularx}
\vspace{0.5em}

\normalsize
Framing from Jaumandreu (2025); the FOC row starts with Doraszelski--Jaumandreu (2013). With a linear $g$ and correct FOCs, all three should agree up to sampling error. Disagreement is informative.
\end{frame}


\begin{frame}{Dynamic Panel (Blundell and Bond 2000)}
With $\omega_{jt} = \rho\,\omega_{jt-1} + \xi_{jt}$, quasi-difference:
\begin{align*}
y_{jt} - \rho\, y_{jt-1} = f(\mathbf{x}_{jt}) - \rho f(\mathbf{x}_{jt-1}) + \xi_{jt} + \varepsilon_{jt} - \rho\,\varepsilon_{jt-1},
\end{align*}
instrumented with $(k_{jt}, k_{jt-1}, l_{jt-1}, m_{jt-1})$ and lagged prices.
\begin{itemize}
\item \cmark No model of input demand: allows unobserved input-price shocks to every input, demand shocks, any conduct, and fixed effects.
\item \xmark Needs $g$ linear; selection is hard to handle; $\varepsilon$ is not separated from $\omega$.
\item \xmark In practice often implausible: in Compustat a Cobb--Douglas dynamic panel gives $\hat\beta_K < 0$ and variable returns above 1.1 (Jaumandreu 2025).
\item For a flexible input, lagged inputs are relevant only if something persistent moves input demand besides $\omega$: firm-specific input prices, or demand differences under market power (Bond et al.\ 2021).
\end{itemize}
\end{frame}


\section{Olley and Pakes (1996)}

\begin{frame}{Olley--Pakes: Deregulating Telecom Equipment}
\begin{itemize}
\item AT\&T's network monopoly extended to equipment: it bought almost everything from its subsidiary Western Electric, and connecting outside equipment was restricted.
\item Two shocks: FCC certification of outside equipment (late 1970s), and the AT\&T divestiture (decreed 1982, effective 1984). Western Electric's share of Bell company purchases fell from 92\% (1982) to 58\% (1986).
\item Massive turnover: plants rose from 133 (1963) to 584 (1987); 79\% of plants active in 1987 entered after 1972, and 60\% of 1972 plants had exited.
\item Questions: did productivity rise, and through which channel, within plants or reallocation across them? Answering needs production function estimates that handle \alert{simultaneity} and \alert{selection}.
\end{itemize}
{\footnotesize Source: Olley and Pakes (1996), Tables I--IV.}
\end{frame}


\begin{frame}{The OP Model}
Incumbents decide each period whether to exit (receiving a scrap value $\underline{V}$), and if not, choose labor and investment:
\begin{align*}
V_t(\omega_t, a_t, k_t) = \max\Big\{\underline{V},\ \sup_{i_t \ge 0}\ \pi_t(\omega_t, a_t, k_t) - c(i_t) + \beta\, \E\big[V_{t+1}(\omega_{t+1}, a_{t+1}, k_{t+1}) \mid \mathcal{I}_t\big]\Big\}
\end{align*}
with $K_{t+1} = (1-\delta_K) K_t + i_t$, $a_{t+1} = a_t + 1$, and $\beta$ the discount factor.
\begin{itemize}
\item Market structure and factor prices are common to all firms, so they are absorbed into the $t$ subscripts (assumed Markov perfect equilibrium).
\item Policies: exit if $\omega_t < \underline{\omega}_t(a_t, k_t)$; invest $i_t = i_t(\omega_t, a_t, k_t)$.
\item \alert{Key result} (Pakes 1994): if $i_t > 0$, investment is strictly increasing in $\omega_t$, so it can be inverted:
\begin{align*}
\omega_{jt} = h_t(i_{jt}, a_{jt}, k_{jt}).
\end{align*}
\item This rests on $\omega$ being the \alert{only} firm-specific unobserved state: no unobserved differences in demand, input prices, or adjustment costs.
\end{itemize}
\end{frame}


\begin{frame}{OP Estimation in Three Steps}
\begin{enumerate}
\item \alert{First stage} (partially linear, Robinson 1988):
\begin{align*}
y_{jt} = \beta_l\, l_{jt} + \Phi_t(i_{jt}, a_{jt}, k_{jt}) + \varepsilon_{jt}, \qquad \Phi_t = \beta_0 + \beta_a a + \beta_k k + h_t(i, a, k).
\end{align*}
Recovers $\hat\beta_l$ and $\hat\Phi_{jt}$ (polynomial or kernel). Capital and age are inside $\Phi$.
\item \alert{Survival}: $\mathcal{P}_{jt} = \Pr(\chi_{jt+1} = 1 \mid i_{jt}, a_{jt}, k_{jt})$, which indexes the exit threshold.
\item \alert{Second stage} (NLLS):
\begin{align*}
y_{jt+1} - \hat\beta_l\, l_{jt+1} = \beta_a a_{jt+1} + \beta_k k_{jt+1} + \tilde g\big(\mathcal{P}_{jt},\ \hat\Phi_{jt} - \beta_a a_{jt} - \beta_k k_{jt}\big) + \xi_{jt+1} + \varepsilon_{jt+1}.
\end{align*}
\end{enumerate}
\begin{itemize}
\item $k_{jt+1}$ was chosen at $t$, so it is uncorrelated with $\xi_{jt+1}$. Labor responds to $\xi_{jt+1}$, which is why $\beta_l$ comes from step 1.
\item Observations with zero investment are dropped (8\% in OP's data; far more elsewhere).
\end{itemize}
\end{frame}


\begin{frame}{OP Results}
\small
\begin{center}
\renewcommand{\arraystretch}{1.15}
\begin{tabular}{lccccc}
\toprule
& \multicolumn{2}{c}{Balanced panel} & \multicolumn{2}{c}{Full sample} & OP \\
\cmidrule(lr){2-3}\cmidrule(lr){4-5}
& OLS & Within & OLS & Within & Kernel \\
\midrule
Labor & .851 & .728 & .693 & .629 & .608 \\
& (.039) & (.049) & (.019) & (.026) & (.027) \\
Capital & .173 & .067 & .304 & .150 & .355 \\
& (.034) & (.049) & (.018) & (.026) & (.058) \\
Obs. & 896 & 896 & 2592 & 2592 & 1758 \\
\bottomrule
\end{tabular}
\end{center}
\begin{itemize}
\item Labor falls and capital rises as simultaneity and selection are corrected, in the directions theory predicts.
\item Within estimators kill the capital coefficient: productivity is not a fixed effect during a restructuring.
\item Overidentification checks (adding $k$ or $a$ inside $\tilde g$) do not reject.
\end{itemize}
{\footnotesize Source: Olley and Pakes (1996), Table VI, columns (1)--(4) and (9).}
\end{frame}


\begin{frame}{The Olley--Pakes Decomposition}
\small
With $\Omega_{jt} = e^{\omega_{jt}}$ and output shares $\lambda_{jt}$, industry productivity is
\begin{align*}
\Omega_t \equiv \sum_j \lambda_{jt}\,\Omega_{jt} = \bar\Omega_t + \sum_j (\lambda_{jt} - \bar\lambda_t)(\Omega_{jt} - \bar\Omega_t),
\end{align*}
the unweighted mean plus a covariance between size and productivity.
\begin{columns}[T]
\begin{column}{0.42\textwidth}
\begin{center}
\scriptsize
\begin{tabular}{lccc}
\toprule
Year & $\Omega_t$ & $\bar\Omega_t$ & Covariance \\
\midrule
1975 & 0.72 & 0.66 & 0.06 \\
1978 & 0.92 & 0.80 & 0.12 \\
1981 & 1.11 & 0.76 & 0.35 \\
1984 & 0.90 & 0.83 & 0.07 \\
1987 & 0.97 & 0.66 & 0.32 \\
\bottomrule
\end{tabular}

{\tiny OP (1996), Table XI.}
\end{center}
\end{column}
\begin{column}{0.56\textwidth}
\begin{itemize}
\item The unweighted mean barely moves; the gains come from output moving toward more productive plants.
\item Across studies the covariance term explains 20--40\% of changes in share-weighted productivity (De Loecker--Syverson 2021).
\item This is revenue productivity.
\item Caveats: share-weighted productivity is not the theoretically correct aggregate (Petrin--Levinsohn 2012); entry and exit need the dynamic version (Melitz--Polanec 2015).
\end{itemize}
\end{column}
\end{columns}
\end{frame}


\section{Levinsohn--Petrin and Ackerberg--Caves--Frazer}

\begin{frame}{Levinsohn and Petrin (2003): Invert Materials Instead}
\small
\begin{columns}[T]
\begin{column}{0.55\textwidth}
\begin{itemize}
\item Use an intermediate input as the proxy: $m_{jt} = \mathbb{M}_t(k_{jt}, \omega_{jt})$, so $\omega_{jt} = h_t(k_{jt}, m_{jt})$.
\item Why:
\begin{itemize}
\item investment is often zero (those observations are dropped) and lumpy (flat spots in $i(\omega)$);
\item materials are static, so monotonicity is a property of $f$ rather than of a dynamic equilibrium;
\item unobserved capital adjustment costs no longer break the inversion.
\end{itemize}
\item Monotonicity from the implicit function theorem on the FOCs for $(m, l)$:
\begin{align*}
\operatorname{sign}\!\Big(\frac{\partial m}{\partial \omega}\Big) = \operatorname{sign}\big(f_{ml} f_{l\omega} - f_{ll} f_{m\omega}\big),
\end{align*}
when the $(m,l)$ Hessian is negative \alert{definite}.
\end{itemize}
\end{column}
\begin{column}{0.43\textwidth}
\begin{center}
\scriptsize
Share of non-zero observations, Chile
\renewcommand{\arraystretch}{1.1}
\begin{tabular}{lcc}
\toprule
Industry & Investment & Materials \\
\midrule
Food (311) & 42.7 & 99.8 \\
Metals (381) & 44.8 & 99.9 \\
Textiles (321) & 41.2 & 99.9 \\
Wood (331) & 35.9 & 99.7 \\
\bottomrule
\end{tabular}

{\tiny LP (2003), Table 2. OP's industry: 8\% zeros.}
\end{center}
\end{column}
\end{columns}
\end{frame}


\begin{frame}{The ACF Critique: Functional Dependence}
LP's first stage, $y_{jt} = \beta_l l_{jt} + \Phi_t(k_{jt}, m_{jt}) + \varepsilon_{jt}$, needs variation in labor that is not explained by $(k, m)$ (Robinson 1988):
\begin{align*}
\E\Big[\big(l_{jt} - \E[l_{jt} \mid k_{jt}, m_{jt}]\big)^2\Big] > 0.
\end{align*}
\begin{itemize}
\item But if labor is chosen like materials, $l_{jt} = \ell_t(k_{jt}, \omega_{jt}) = \ell_t\big(k_{jt}, h_t(k_{jt}, m_{jt})\big)$: \alert{no variation left}.
\item The simultaneity algebra again, now in LP's first stage. Cobb--Douglas with a price-taking FOC for $m$, inverted for $\omega$ and plugged back in:
\begin{align*}
y_{jt} = \ln\Big(\frac{1}{\beta_m}\Big) + \ln\Big(\frac{w^M_t}{p_t}\Big) + m_{jt} + \varepsilon_{jt}.
\end{align*}
$\beta_l$ does not appear at all.
\item With polynomials you will still get a number. It just is not a consistent estimate of $\beta_l$.
\item The same argument applies to OP's investment proxy.
\end{itemize}
\end{frame}


\begin{frame}{What Would Break the Functional Dependence?}
ACF search for data-generating processes consistent with LP that still identify $\beta_l$ in the first stage. Only knife-edge cases work:
\begin{enumerate}
\item iid \alert{optimization error} in labor, but not in materials (workers calling in sick, with materials ordered beforehand);
\item an iid shock to wages or demand realized \alert{after} $m$ is chosen and \alert{before} $l$;
\item (OP only) labor chosen at $t-b$, on $\omega_{jt-b}$, and non-dynamic, while investment is chosen at $t$.
\end{enumerate}
What does \alert{not} help:
\begin{itemize}
\item observed firm-level input prices (they enter $\Phi_t$ too);
\item dynamic labor ($l_{jt-1}$ becomes a state in both input demands);
\item classical measurement error.
\end{itemize}
\end{frame}


\begin{frame}{The ACF Procedure}
\small
Invert the demand for materials \alert{conditional on labor}, $m_{jt} = \tilde{\mathbb{M}}_t(k_{jt}, l_{jt}, \omega_{jt})$, so $\omega_{jt} = h_t(k_{jt}, l_{jt}, m_{jt})$.
\begin{enumerate}
\item First stage identifies \alert{nothing} but nets out $\varepsilon$: \quad $y_{jt} = \Phi_t(k_{jt}, l_{jt}, m_{jt}) + \varepsilon_{jt}$.
\item For trial $\beta$: $\omega_{jt}(\beta) = \hat\Phi_{jt} - \beta_k k_{jt} - \beta_l l_{jt}$; regress on $\omega_{jt-1}(\beta)$ to get $\xi_{jt}(\beta)$; solve
\begin{align*}
\E\Big[\xi_{jt}(\beta) \otimes \big(1,\ k_{jt},\ l_{jt-1},\ \hat\Phi_{jt-1}\big)'\Big] = 0.
\end{align*}
\end{enumerate}
\begin{columns}[T]
\begin{column}{0.48\textwidth}
\begin{itemize}
\item Allows dynamic labor, labor chosen before materials, and serially correlated unobserved wage shocks.
\item Does \alert{not} allow unobserved materials prices or demand shocks.
\item Meant for value added, or gross output that is Leontief in $m$. Next: why not general gross output.
\item Wooldridge (2009): both stages as one GMM system; efficient, easy robust SEs.
\end{itemize}
\end{column}
\begin{column}{0.5\textwidth}
\begin{center}
\scriptsize
Serially correlated unobserved cost shocks allowed?
\renewcommand{\arraystretch}{1.1}
\begin{tabular}{lccc}
\toprule
 & to $k$ & to $l$ & to $m$ \\
\midrule
OP & \xmark & \xmark & \xmark \\
LP & \cmark & \xmark & \xmark \\
ACF & \cmark & \cmark & \xmark \\
Dynamic panel & \cmark & \cmark & \cmark \\
\bottomrule
\end{tabular}
\end{center}
\end{column}
\end{columns}
\end{frame}


\begin{frame}{ACF Monte Carlo}
\small
Leontief in materials; true $(\beta_l, \beta_k) = (0.6, 0.4)$; 1000 firms $\times$ 10 periods; no measurement error in $m$.
\begin{center}
\renewcommand{\arraystretch}{1.15}
\begin{tabular}{lcccc}
\toprule
& \multicolumn{2}{c}{ACF} & \multicolumn{2}{c}{LP} \\
\cmidrule(lr){2-3}\cmidrule(lr){4-5}
DGP & $\hat\beta_l$ & $\hat\beta_k$ & $\hat\beta_l$ & $\hat\beta_k$ \\
\midrule
1: serially correlated wages, $l$ chosen at $t-b$ & 0.600 & 0.399 & 0.000 & 1.121 \\
2: optimization error in $l$ (LP's case) & 0.600 & 0.400 & 0.600 & 0.399 \\
3: both & 0.596 & 0.406 & 0.473 & 0.588 \\
\bottomrule
\end{tabular}
\end{center}
\begin{itemize}
\item DGP 1: functional dependence in action, as LP's labor coefficient collapses to zero.
\item DGP 2: LP is consistent and more precise; ACF pays an efficiency cost for robustness.
\item OLS gives $\hat\beta_l \approx 0.85$--$0.95$. OP with firm-specific capital adjustment costs gives $\hat\beta_l \approx 0.82$.
\item Adding measurement error to $m$ moves LP much more than ACF.
\end{itemize}
{\footnotesize Source: Ackerberg, Caves and Frazer (2015), Tables I--II.}
\end{frame}


\section{Gross Output}

\begin{frame}{Gross Output Is Not Identified by These Moments (GNR 2020)}
Put materials back in: $y_{jt} = f(k_{jt}, l_{jt}, m_{jt}) + \omega_{jt} + \varepsilon_{jt}$, with $m$ flexible. After conditioning on the instruments, what still moves materials?
\begin{align*}
m_{jt} = \mathbb{M}\big(k_{jt},\ l_{jt},\ g(h(\mathbf{o}_{jt-1})) + \xi_{jt} - d_t\big), \qquad d_t = \ln(w^M_t / p_t).
\end{align*}
Only $\xi_{jt}$, which is orthogonal to every instrument by construction, and the aggregate price ratio $d_t$.
\begin{itemize}
\item \alert{Theorem 1}: with $d_t = d$, for every $a \in (0,1)$
\begin{align*}
\tilde f = (1-a)\, f + a\, \Phi, \qquad \tilde g(x) = a d + (1-a)\, g\Big(\frac{x - a d}{1 - a}\Big)
\end{align*}
satisfies exactly the same moment conditions. The same holds for investment proxies and for dynamic panel (their App. O1).
\item Time-series price variation identifies in theory but not in practice. In their Monte Carlo the estimated $\theta^M$ is about 0.9--1.0 against a truth of 0.65, unless price variation is 10 times the baseline and $T \ge 20$.
\item What you need is an \alert{excluded shifter} of materials demand that varies across firms ($\sigma^2_d > 0$: firm-level input prices), or known returns to scale instead (FGT 2019).
\end{itemize}
\end{frame}


\begin{frame}{GNR's Fix: Use the First-Order Condition}
Under price-taking in output and materials, take logs of the FOC $p_t\, \partial F/\partial M\, e^{\omega}\, \mathcal{E} = w^M_t$ and subtract the production function:
\begin{align*}
\ln \alpha^M_{jt} = \ln \mathcal{E} + \ln \theta^M(k_{jt}, l_{jt}, m_{jt}) - \varepsilon_{jt}, \qquad \mathcal{E} = \E[e^{\varepsilon}].
\end{align*}
\begin{enumerate}
\item A nonparametric \alert{share regression} identifies $\theta^M(\cdot)$ and $\varepsilon_{jt}$.
\item Integrate the elasticity: $\int \theta^M\, dm = f(k, l, m) + \mathcal{C}(k, l)$.
\item The remainder obeys a dynamic panel structure:
\begin{align*}
\mathcal{Y}_{jt} \equiv y_{jt} - \varepsilon_{jt} - \int \theta^M dm_{jt} = -\mathcal{C}(k_{jt}, l_{jt}) + \omega_{jt},
\end{align*}
with moments $\E[\xi_{jt} \mid k_{jt}, l_{jt}, \mathcal{Y}_{jt-1}, k_{jt-1}, l_{jt-1}] = 0$.
\end{enumerate}
A nonparametric version of Griliches--Ringstad (1971): index-number logic for the flexible input, lags for the rest. Estimation uses polynomial sieves with closed-form integrals and a bootstrap.
\end{frame}


\begin{frame}{GNR Estimates: Chile and Colombia}
\small
\begin{center}
\renewcommand{\arraystretch}{1.15}
\begin{tabular}{lcccccc}
\toprule
& \multicolumn{2}{c}{Colombia, all mfg.} & \multicolumn{2}{c}{Chile, all mfg.} & \multicolumn{2}{c}{Chile, food (311)} \\
\cmidrule(lr){2-3}\cmidrule(lr){4-5}\cmidrule(lr){6-7}
Elasticity & GNR & OLS & GNR & OLS & GNR & OLS \\
\midrule
Labor & .35 & .26 & .38 & .20 & .28 & .17 \\
Capital & .14 & .06 & .16 & .09 & .11 & .05 \\
Intermediates & .54 & .72 & .55 & .77 & .67 & .83 \\
Sum & 1.04 & 1.04 & 1.09 & 1.06 & 1.05 & 1.05 \\
\bottomrule
\end{tabular}
\end{center}
\begin{itemize}
\item OLS overstates the intermediate-input elasticity (by 34\% on average) and understates labor and capital.
\item OLS also understates productivity dispersion (Colombia: 95/5 ratio of 2.36 under GNR vs 1.94 under OLS) and persistence (0.64 vs 0.53).
\item Robust to fixed effects, to treating part of intermediates as non-flexible, and to large adjustment costs in materials (Monte Carlo).
\end{itemize}
{\footnotesize Source: Gandhi, Navarro and Rivers (2020), Tables 2--3.}
\end{frame}


\begin{frame}{Leontief in Materials: Why It Is Useful}
\footnotesize
\textit{Structural value added}: \quad $Y_{jt} = \min\big\{F(K_{jt}, L_{jt})\, e^{\omega_{jt}},\ \gamma M_{jt}\big\}\, e^{\varepsilon_{jt}}$; a producing firm sets $\gamma M = F(K,L)\,e^{\omega}$.
\begin{itemize}
\item \alert{Identification}: materials drop out of what you estimate, $y = f(k,l) + \omega + \varepsilon$. Yet they remain a perfect proxy, $m = f(k,l) + \omega - \ln\gamma$, monotone in $\omega$ given $(k,l)$. GNR's non-identification result has no flexible input to bite on. This is the case ACF have in mind.
\item \alert{Plausible} at a disaggregated level, where materials are a recipe: steel mills (Hendel--Spiegel 2014), brewing (De Loecker--Scott 2016), oil wells (Asker--Collard-Wexler--De Loecker 2019), and cigarettes, with 686 mg of leaf per 1000 mg cigarette and s.d.\ 30 mg across Chinese manufacturers (Rubens 2023).
\item \alert{Markups}: at the kink materials have no FOC, so use labor. But marginal cost includes materials:
\begin{align*}
mc = \frac{w^M}{\gamma} + \frac{w^L L}{\theta^L Y} \quad\Longrightarrow\quad \mu = \Big(\alpha^M + \frac{\alpha^L}{\theta^L}\Big)^{-1}.
\end{align*}
The naive $\theta^L/\alpha^L$ overstates $\mu$: with $\alpha^M = 0.6$, $\alpha^L = 0.2$, $\theta^L = 0.6$, it gives 3.0 against the true 1.07 (De Loecker--Scott; Rubens 2023, eq.\ 4d).
\item \alert{Cautions}: idle capacity breaks the Leontief equality (GNR's critique, ACF App.); estimated labor--materials substitution elasticities are well above zero (Doraszelski--Jaumandreu 2018; Edmond--Midrigan--Xu calibrate 0.5).
\end{itemize}
\end{frame}


\section{Endogenous Productivity}

\begin{frame}{Endogenous Productivity: Firm Actions Belong in $g$}
\small
Exporting, R\&D, or trade protection at $t-1$ shift the conditional mean of productivity: $\omega_{jt} = g(\omega_{jt-1}, a_{jt-1}) + \xi_{jt}$, with $a_{jt-1} \in \mathcal{I}_{jt-1}$. The moments are unchanged; $g$ just gains an argument.

Suppose the truth is $\omega_{jt} = \omega_{jt-1} + \gamma\, e_{jt-1} + \xi_{jt}$ with export status $e$, but $g$ is treated as exogenous (De Loecker 2013). With $\tilde y = y - \hat\beta_l l$:
\begin{align*}
\Delta\tilde y_{jt} = \beta_k \Delta k_{jt} + \underbrace{\gamma\, e_{jt-1} + \xi_{jt}}_{\text{error}} + \Delta\varepsilon_{jt}
\quad\Longrightarrow\quad
\operatorname{plim} \hat\beta_k = \beta_k + \gamma\,\frac{\Cov(\Delta k_{jt}, e_{jt-1})}{\Var(\Delta k_{jt})} > \beta_k.
\end{align*}
\begin{itemize}
\item Exporters invest more, so $\hat\beta_k$ is too high, and a DID on $\hat\omega$ subtracts too much of export starters' capital growth. Slovenia: $\hat\beta_k$ is about 30\% higher under an exogenous $g$, and the DID misses 1.5 pp of a 4.1\% learning-by-exporting effect in the entry year.
\item Estimating $\hat\omega$ with an exogenous $g$ and then regressing it on $e$ tests what step 1 assumed away.
\item \alert{R\&D} (Doraszelski--Jaumandreu 2013, Spain): the exogenous process is rejected in all 9 industries; R\&D makes productivity riskier (larger $\hat\xi^2$ in every industry, while physical investment does not); the net return is about 40\%, roughly twice that on physical capital.
\end{itemize}
\end{frame}


\section{Markups from Production Data}

\begin{frame}{Three Ways to Measure a Markup}
$\mu_{jt} \equiv p_{jt}/mc_{jt}$. Marginal cost is rarely observed, so something has to stand in for it.
\begin{itemize}
\item \alert{Accounting}: revenue over total cost. Right only if average cost equals marginal cost (CRS, no fixed costs).
\item \alert{Demand approach} (NEIO; Weeks 2--7): estimate demand elasticities, assume a conduct model, and back out $mc$ from the pricing FOC.
\begin{itemize}
\item needs product-level prices and quantities, instruments, and a model of competition;
\item delivers counterfactuals (mergers), which the production approach cannot.
\end{itemize}
\item \alert{Production approach} (Hall 1988; De Loecker--Warzynski 2012): cost minimization in one frictionless variable input.
\begin{itemize}
\item no demand system and no conduct assumption;
\item needs an output elasticity, hence a production function;
\item works across many industries at once, so it has driven the macro debate on market power.
\end{itemize}
\item They rest on different assumptions, so agreement is informative: De Loecker and Scott (2016) find the two approaches agree for US brewing.
\end{itemize}
\end{frame}


\begin{frame}{De Loecker and Warzynski (2012): The Markup Formula}
Minimize cost given dynamic inputs (capital) and productivity:
\begin{align*}
\min_{\{X\}_{X \in \mathcal{V}}} \sum_{X \in \mathcal{V}} w^X X \quad \text{s.t.}\quad F(\mathbf{X}, K)\, e^{\omega} \ge Y
\quad\Longrightarrow\quad w^X = mc \cdot \frac{\partial Y}{\partial X}.
\end{align*}
Multiply by $X/Y$, then multiply and divide by $p$:
\begin{align*}
\theta^X_{jt} = \frac{p_{jt}}{mc_{jt}} \cdot \frac{w^X_{jt} X_{jt}}{p_{jt} Y_{jt}} = \mu_{jt}\, \alpha^X_{jt}
\qquad\Longrightarrow\qquad
\boxed{\mu_{jt} = \frac{\theta^X_{jt}}{\alpha^X_{jt}}}
\end{align*}
\begin{itemize}
\item Requires: cost minimization; $X$ free of adjustment costs; firm is a price taker in the market for $X$; prices set period by period.
\item Does \alert{not} require: a demand system, a conduct model, constant returns, or the user cost of capital (we condition on $K$).
\item The formula is old (Hall 1988 assumed $\mu$ constant across firms). What is new is firm-level $\theta^X_{jt}$ from a carefully estimated production function.
\end{itemize}
\end{frame}


\begin{frame}{Implementing DLW}
\begin{itemize}
\item Estimate $f$ by ACF, with extra input-demand shifters $\mathbf{d}_{jt}$ in the proxy: $m_{jt} = \mathbb{M}_t(k_{jt}, \omega_{jt}, \mathbf{d}_{jt})$.
\begin{itemize}
\item DLW include input prices and export status. De Loecker et al.\ (2016) add output prices, market shares, product dummies, and tariffs.
\item This is the only restriction on competition: \alert{everything} that shifts input demand across firms must be in $\mathbf{d}$.
\item Monotonicity under imperfect competition: more productive firms must not set ``inordinately'' higher markups (Melitz--Levinsohn).
\end{itemize}
\item Translog: $\hat\theta^L_{jt} = \hat\beta_l + 2\hat\beta_{ll}\, l_{jt} + \hat\beta_{lk}\, k_{jt}$. Under Cobb--Douglas $\theta$ is constant, so \alert{all} variation in the share becomes variation in the markup.
\item Use planned rather than realized output in the share: $\hat\alpha^X_{jt} = w^X_{jt} X_{jt} / \big(p_{jt} Y_{jt}\, e^{-\hat\varepsilon_{jt}}\big)$.
\item Then project: $\ln \hat\mu_{jt} = \delta_0 + \delta_1\, \text{export}_{jt} + \text{controls} + e_{jt}$.
\end{itemize}
\end{frame}


\begin{frame}{Multi-Product Firms (De Loecker, Goldberg, Khandelwal, Pavcnik 2016)}
\small
Indian manufacturers, 1989--2003: product quantities and unit-value prices, but inputs only at the firm level, as deflated expenditures $\tilde{\mathbf{x}}_{jt}$. Products $n$, input shares $\phi_{jnt}$.
\begin{enumerate}
\item Estimate $f$ on single-product firms (selection-corrected), proxying input prices by $\ln\tilde w = \mathcal{W}(\ln p, \text{market share}, \text{product})$.
\item $\omega_{jt}$ is common to the firm's products, so physical output pins down each product's input share:
\[ \hat y_{jnt} = f\big(\tilde{\mathbf{x}}_{jt} + \ln\phi_{jnt} - \ln\hat{\tilde w}_{jnt}\big) + \omega_{jt}, \qquad \textstyle\sum_n \phi_{jnt} = 1. \]
Cobb--Douglas with returns $\Gamma$: $\phi_{jnt} \propto \hat{\tilde w}_{jnt}\,\hat Y_{jnt}^{1/\Gamma}$. Splitting inputs by revenue is right only if markups are equal across products.
\item $\mu_{jnt} = \theta^M_{jnt}\, p_{jnt}Y_{jnt}/(\phi_{jnt}\, w^M_{jt}M_{jt})$, and $mc_{jnt} = p_{jnt}/\mu_{jnt}$.
\end{enumerate}
\begin{columns}[T]
\begin{column}{0.55\textwidth}
\centering
\scriptsize
\begin{tabular}{lccc}
\toprule
India 1989--97, \% & output tariffs & input tariffs & total \\
\midrule
$\Delta\ln p$ & $-9.7$ & $-8.4$ & $-18.1$ \\
$\Delta\ln mc$ & $-2.9$ & $-27.8$ & $-30.7$ \\
$\Delta\ln\mu$ & $-6.8$ & $+19.4$ & $+12.6$ \\
\bottomrule
\end{tabular}
\end{column}
\begin{column}{0.43\textwidth}
\scriptsize
Tariffs fell 62 pp (output) and 24 pp (input). \alert{Cheaper inputs were kept as markup}: pass-through about 0.3. Median markup 1.34 (mean 2.70). The markup total is not significant (s.e.\ 11.9). Closed form and split derived here (eqs.\ 23--24, App.\ B, Table IX).
\end{column}
\end{columns}
\end{frame}


\begin{frame}{One Markup, Many Inputs}
\small
If the FOC holds for every variable input, $\theta^X = \mu\, \alpha^X$ for all $X \in \mathcal{V}$. Sum over $X$, with $\nu_{jt} \equiv \sum_{X \in \mathcal{V}} \theta^X_{jt}$ the short-run scale elasticity ($\nu$ here is not a random-coefficient draw):
\begin{align*}
\nu_{jt} = \mu_{jt}\, \alpha^{\mathcal{V}}_{jt}
\qquad\Longrightarrow\qquad
\boxed{\mu_{jt} = \nu_{jt}\, \frac{p_{jt} Y_{jt}}{VC_{jt}}},
\qquad \frac{\theta^X_{jt}}{\nu_{jt}} = \frac{w^X_{jt} X_{jt}}{VC_{jt}}.
\end{align*}
\begin{itemize}
\item The markup is short-run returns to scale times revenue over variable cost. Equivalently, $mc = AVC/\nu$ (Jaumandreu 2025). With CRS and all inputs variable, this is the accounting markup.
\item What you really need is $\nu$, not each $\theta^X$.
\item With two flexible inputs the markup is \alert{overidentified}: labor and materials give the same $\mu$ iff
\begin{align*}
\frac{\theta^L_{jt}}{\theta^M_{jt}} = \frac{w^L_{jt} L_{jt}}{w^M_{jt} M_{jt}}.
\end{align*}
This is exactly Raval's (2023) test, and it is what fails in the data.
\end{itemize}
{\footnotesize Combining DLW's formula across inputs; derived here.}
\end{frame}


\begin{frame}{DLW Results: Slovenian Manufacturing, 1994--2000}
\small
\begin{columns}[T]
\begin{column}{0.44\textwidth}
\begin{center}
\scriptsize
Median markups\par\smallskip
\renewcommand{\arraystretch}{1.1}
\begin{tabular}{lc}
\toprule
Hall (1986) & 1.03 \\
Klette (1999) & 1.12 \\
Cobb--Douglas, value added & 1.17 \\
Translog, value added & 1.28 \\
Gross output, labor FOC & 1.26 \\
Gross output, materials FOC & 1.22 \\
Single markup, levels & 1.16 \\
Single markup, first differences & 1.11 \\
\bottomrule
\end{tabular}

{\tiny DLW Table 2; s.d.\ of firm markups about 0.5.}
\end{center}
\end{column}
\begin{column}{0.54\textwidth}
\begin{itemize}
\item 7,915 firms. Controlling for productivity raises the estimates relative to Hall and Klette.
\item Using capital's ``FOC'' gives 1.5--1.6. That wedge includes adjustment costs, which is why capital is not used.
\item Exporters charge higher markups: about 8\% (0.13--0.18 in levels). This falls to 2\% once productivity is controlled, so most of the gap is lower marginal cost.
\item Markups rise when firms start exporting.
\end{itemize}
\end{column}
\end{columns}
\end{frame}


\begin{frame}{One FOC, Two Unknowns}
\small
The flexible-input FOC, $\ln \alpha^M_{jt} = \ln \theta^M_{jt} - \ln \mu_{jt} + \ln \mathcal{E} - \varepsilon_{jt}$, is one equation in two unknowns.
\begin{itemize}
\item \alert{GNR} read it left to right with $\mu \equiv 1$ (price taking): the share regression identifies $\theta^M$. Under market power it recovers $\theta^M/\mu$ instead.
\item \alert{DLW} read it right to left: take $\theta^M$ from the production function and back out $\mu$. The markup is identified by whatever identifies $\theta^M$: an \alert{excluded shifter} of input demand, here serially correlated firm-level input prices (DLW p.\ 2446). Unobserved, those prices would also break the proxy (Flynn--Gandhi--Traina 2019).
\item \alert{Flynn--Gandhi--Traina (2019)} use a shape restriction instead. With Cobb--Douglas and $\pi_x \equiv \partial\E[x_{jt} \mid k_{jt}, \mathbf{o}_{jt-1}]/\partial k_{jt}$, timing gives one equation, $\pi_y = \theta^K + \theta^M \pi_m$. Constant returns close it:
\begin{align*}
\theta^M = \frac{1 - \pi_y}{1 - \pi_m}, \qquad \mu = (1 - \theta^K)\,\frac{pY}{VC}.
\end{align*}
This fails when $\pi_m = 1$, which (with a constant markup) happens exactly when $\mu = 1$: GNR's case. It also needs \alert{quantity} output. With revenue and a constant markup, planned revenue is $\ln(\mu/\theta^M) + \ln w^M + m$, so $\pi_r = \pi_m$ and the same moments return $\hat\theta^M = 1$ and $\hat\mu = pY/VC = \mu/\theta^M$: Bond et al.'s problem in CRS form.
\end{itemize}
{\footnotesize Derived here from GNR (2020), DLW (2012) and FGT (2019).}
\end{frame}


\section{Where the Production Approach Breaks}

\begin{frame}{Estimate the Demand Side Too (De Loecker 2011)}
\small
Belgian textiles, 1994--2002: no prices, but segment output $\bar y_{st}$ (including imports), EU quota protection $\mathit{qr}_{jt} \in [0,1]$, and product dummies $\mathbf{D}_j$. Take the TFPR slide's equation with $\eta \to \eta_s$ and $\delta_{jt} = \mathbf{D}_j'\boldsymbol{\kappa} + \tau\,\mathit{qr}_{jt} + \tilde\delta_{jt}$:
\begin{align*}
\tilde r_{jt} &= \textstyle\sum_X \underbrace{\tfrac{\theta^X}{\mu_s}}_{\beta_X} x_{jt} + \underbrace{\tfrac{1}{\eta_s}}_{\beta_s} \bar y_{st} + \tfrac{1}{\eta_s}\big(\tau\,\mathit{qr}_{jt} + \mathbf{D}_j'\boldsymbol{\kappa}\big) + \tfrac{\eta_s - 1}{\eta_s}\,\omega_{jt} + \tilde\varepsilon_{jt}, \\
\omega_{jt} &= g(\omega_{jt-1}, \mathit{qr}_{jt-1}) + \xi_{jt}, \qquad \omega_{jt} = h_t(k_{jt}, m_{jt}, \mathit{qr}_{jt}, \bar y_{st}, \mathbf{D}_j).
\end{align*}
\begin{itemize}
\item The coefficient on market size is the Lerner index: $\hat\mu_s = 1/(1 - \hat\beta_s)$, $\hat\theta^X = \hat\mu_s\,\hat\beta_X$. Quotas shift demand at once ($\tau$), productivity only with a lag ($g$).
\item With demand controls: $\hat\eta$ 2.86 $\to$ 3.85, $\hat\mu$ 1.54 $\to$ 1.35, returns to scale 1.3 $\to$ 1.11 (revenue OLS: 0.95).
\item Removing all quotas raises productivity 2.1\%, not 8.6\% (over three years: 4.6\%, not 13.5\%).
\item CES fixes $\mu_s$, so liberalization cannot move markups; $\beta_X$ is the next slide's revenue elasticity.
\end{itemize}
\end{frame}


\begin{frame}{Revenue Elasticities Carry No Markup Information (Bond et al.\ 2021)}
With revenue data we can at best learn the \alert{revenue} elasticity. With $R = pY$:
\begin{align*}
\theta^X_R \equiv \frac{\partial r_{jt}}{\partial x_{jt}} = \frac{\partial \ln R}{\partial \ln Y}\,\theta^X_{jt} = \frac{MR_{jt}}{p_{jt}}\,\theta^X_{jt}
\quad\Longrightarrow\quad
\frac{\theta^X_R}{\alpha^X_{jt}} = \frac{MR_{jt}}{p_{jt}}\,\mu_{jt} = \frac{MR_{jt}}{mc_{jt}} = 1.
\end{align*}
\begin{itemize}
\item This holds under profit maximization for \alert{any} demand system and market structure (their App.\ A.1: Bertrand and Cournot).
\item Klette--Griliches said revenue biases elasticities down. This is the strong form: the firm's own revenue elasticity contains no markup information.
\item Revenue data cannot even identify $\theta^X$ nonparametrically. $m$ and $p$ solve the same problem, $m = m^*(k,l,\omega,\delta,w^M)$ and $p = p^*(k,l,\omega,\delta,w^M)$, which is ACF-style functional dependence.
\item With quantity data, identification needs persistent firm differences in $w^M$ \emph{or} in demand. Market power turns demand heterogeneity into instrument relevance.
\item Their bottom line: \alert{levels need firm-level output prices}. Mean differences between groups with a common technology (exporters) survive; trends need a stable $\theta$.
\end{itemize}
\end{frame}


\begin{frame}{\ldots But Dispersion and Trends Survive (De Ridder, Grassi, Morzenti 2026)}
\footnotesize
Estimators never deliver firm $j$'s own revenue elasticity; they deliver a weighted average across firms. Cobb--Douglas, log variable input $v_{jt}$, lagged input as instrument:
\begin{align*}
\hat\theta = \theta + \frac{\E[\ln p_{jt}\, v_{jt-1}]}{\E[v_{jt}\, v_{jt-1}]},
\qquad
\hat\mu^R_{jt} = \underbrace{\E\!\left[\frac{\varpi_{jt}}{\mu_{jt}}\right]}_{\text{common constant}} \mu_{jt},
\qquad \varpi_{jt} = \frac{v_{jt}\, v_{jt-1}}{\E[v_{jt}\, v_{jt-1}]}.
\end{align*}
\begin{itemize}
\item So $\operatorname{corr}(\ln\hat\mu^R, \ln\mu) = 1$ while the \alert{level} is uninformative (its mean is about 1). With identical demand elasticities (CES, atomistic firms) every firm has $\hat\mu^R = 1$: the Bond et al.\ case.
\item Translog: $\hat\mu^R_{jt} = (1 + f(v_{jt}))\,\mu_{jt}$. The correlation stays high when prices are close to orthogonal to the instruments: 0.94 in an Atkeson--Burstein simulation.
\item France 2009--19, with firm prices from unit values: corr(revenue markup, quantity markup) = 0.63 within sectors (0.80 in first differences). Mean markup is 1.08 from revenue vs 1.45 from quantities, but 2010--19 trends match, and the aggregation method moves the level, not the trend.
\item \alert{Rule}: levels need prices. Dispersion and trends survive revenue data if technology is common across firms (and stable over time).
\end{itemize}
\end{frame}


\begin{frame}{When Productivity Cannot Be Inverted (Doraszelski and Li 2025)}
\small
The proxy assumes \alert{invertibility}: $\E[\omega_{jt} \mid \mathbf{o}_{jt}] = \omega_{jt}$. It fails with demand shocks $\delta_{jt}$, unobserved rivals, changes in conduct (ownership), or unobserved input prices.
\begin{itemize}
\item \alert{Test}: invertibility implies $\E[y_{jt} \mid \mathbf{o}_{jt}, \mathbf{o}_{jt-1}] = \E[y_{jt} \mid \mathbf{o}_{jt}]$. Estimate $y_{jt} = \mathbf{o}_{jt-1}'\beta + \psi(\mathbf{o}_{jt}) + e_{jt}$ and test $\beta = 0$. Rejected almost everywhere (Spain ESEE; US NBER-CES).
\item \alert{Consequence}: step 1 delivers planned output minus $\zeta_{jt} = \omega_{jt} - \E[\omega_{jt} \mid \mathbf{o}_{jt}]$. With an AR(1), the step-2 moment is off by $\rho\, \E[\zeta_{jt-1} \mid \mathbf{z}_{jt}]$, and \alert{$k_{jt}$ is no longer a valid instrument} (it was chosen knowing $\omega_{jt-1}$).
\item \alert{Fix}: make the step-1 conditioning set contain the instruments, $\mathbf{o}_{jt-1} \supseteq \mathbf{z}_{jt}$, so add the lead $k_{jt+1}$ to step 1. This is exact for an AR(1) and a first-order correction otherwise. Add every demand and conduct covariate you have; a Neyman-orthogonal step 2 lets step 1 use ML.
\end{itemize}
\begin{center}
\scriptsize
\renewcommand{\arraystretch}{1.1}
\begin{tabular}{lccc}
\toprule
Monte Carlo (truth: mean log markup 0.25, $g'(0) = 0.70$) & Standard ACF & $+\,k_{t+1}$ & $+\,k_{t+1}, p_t$ \\
\midrule
AR(1) productivity: mean log markup & $-0.24$ & 0.25 & 0.25 \\
Nonlinear productivity: mean log markup & $-0.15$ & 0.30 & 0.25 \\
\bottomrule
\end{tabular}
\end{center}
\end{frame}


\begin{frame}{Marginal Cost as a Sufficient Statistic (Jaumandreu 2025)}
\begin{itemize}
\item With market power the FOC is $MR_{jt}\, \partial F/\partial X\, e^{\omega_{jt}} = w^X_{jt}$. Marginal revenue depends on conduct, rivals, and unobserved demand: it is a correspondence, not a function of observables. Proxies built on $p_{jt}$ (including $p$ in the ACF first stage) are inconsistent.
\item But $MR = mc$, and $mc$ is a single-valued function of $(k, w, Y, \omega)$: a \alert{sufficient statistic} for conduct and demand. With no FOC wedges, $mc = AVC/\nu$.
\item Cost-share proxy, robust to any conduct (needs $g$ linear):
\begin{align*}
\omega_{jt-1} = y_{jt-1} - f(\mathbf{x}_{jt-1}) + \ln\frac{\nu_{jt-1}\, \alpha^X_{jt-1}/\alpha^{\mathcal{V}}_{jt-1}}{\theta^X_{jt-1}} - \varepsilon_{jt-1}.
\end{align*}
\item Two Hausman tests: market power (cost-share vs revenue-share proxy) and specification (dynamic panel vs cost-share proxy).
\item Compustat 1960--2018 (5,621 firms): perfect competition is rejected; Cobb--Douglas fails the specification test (dynamic panel $\hat\beta_K < 0$, $\hat\nu > 1.1$); a translog with labor-augmenting productivity passes, with $\hat\nu = 0.70$--$0.85$.
\end{itemize}
\end{frame}


\begin{frame}{Testing the Production Approach (Raval 2023)}
\small
Any flexible input should give the same markup. ACF, translog, Hicks-neutral productivity, seven datasets:
\begin{center}
\scriptsize
\renewcommand{\arraystretch}{1.1}
\begin{tabular}{lcccc}
\toprule
& \multicolumn{3}{c}{90/50 ratio of markups} & Elasticity of $\mu^L$ \\
\cmidrule(lr){2-4}
Dataset & Labor & Materials & Labor + materials & w.r.t.\ $\mu^M$ \\
\midrule
Chile & 2.03 & 1.39 & 1.17 & $-0.16$ \\
Colombia & 1.82 & 1.43 & 1.17 & $-0.28$ \\
India & 2.95 & 1.29 & 1.14 & $-0.53$ \\
Indonesia & 3.12 & 1.46 & 1.16 & $-0.48$ \\
US (Compustat) & 3.44 & 2.25 & 2.13 & $-0.50$ \\
Southern Europe & 1.79 & 1.27 & 1.10 & $-0.32$ \\
US retailer (stores) & 1.30 & 1.03 & 1.02 & $-10.08$ \\
\bottomrule
\end{tabular}

{\tiny Raval (2023), Tables 2--3; regressions include year and industry effects.}
\end{center}
\begin{itemize}
\item Labor markups are far more dispersed, and labor and materials markups are \alert{negatively} correlated everywhere.
\item Not labor frictions: energy and raw-materials markups also disagree. Not local input prices or monopsony: the retailer result survives MSA-by-year effects.
\end{itemize}
\end{frame}


\begin{frame}{Opposite Trends (Raval 2023, US Compustat)}
\begin{center}
\includegraphics[width=0.9\textwidth]{raval_fig3_us_trends.png}
\end{center}
{\small Markup change relative to 1970, translog: (a) unweighted, (b) cost-weighted. Labor (red) +91\% and +154\%; materials (blue) $-46\%$ and $-23\%$; labor + materials (green) +3\%. Which input you choose decides whether US markups rose.}
\end{frame}


\begin{frame}{Why: Labor-Augmenting Productivity}
\small
CES with Hicks-neutral $\omega$ and labor-augmenting $\omega^L$ (Raval 2023, eqs.\ 8--10):
\begin{align*}
\frac{w^L_{jt} L_{jt}}{w^M_{jt} M_{jt}} = \Big(\frac{w^L_{jt}}{w^M_{jt}}\Big)^{1-\sigma}\Big(\frac{\alpha_l}{\alpha_m}\Big)^{\sigma} e^{(\sigma - 1)\,\omega^L_{jt}},
\end{align*}
where $\sigma$ is the elasticity of substitution (not a choice probability).
\begin{itemize}
\item With $\sigma < 1$, higher $\omega^L$ \alert{lowers} labor's cost share. A Hicks-neutral estimate attributes this to markups: $\hat\mu^L$ too high, $\hat\mu^M$ too low, hence the negative correlation.
\item Cobb--Douglas makes $\theta^L/\theta^M$ a constant, and translog a deterministic function of inputs. Neither has an error term to absorb variation in cost ratios.
\item Evidence that $\sigma < 1$: Raval (2019, RAND) estimates a capital--labor $\sigma \approx 0.3$--$0.5$ from wage differences across 722 US commuting zones (amenity, Bartik, and BGS instruments). Labor-augmenting productivity is persistent and correlated with size, exports, and growth.
\item Doraszelski and Jaumandreu (2019): DLW's exporter premium \alert{flips sign} when computed from materials instead of labor (Spain; Slovenia).
\end{itemize}
\end{frame}


\begin{frame}{Allowing for Heterogeneous Technology}
\small
\begin{itemize}
\item \alert{Demirer (2026)}: gross output with Hicks-neutral and labor-augmenting productivity; only homothetic separability of $(L, M)$.
\begin{itemize}
\item Invert the ratio of flexible-input demands for $\omega^L$, then use an Imbens--Newey control variable for $\omega$.
\item Functional dependence again: input--output variation identifies only the \alert{sum} $\nu = \theta^L + \theta^M$ (exactly what the markup needs). The split comes from the FOCs, $\theta^L/\theta^M = \alpha^L/\alpha^M$, so labor and materials markups agree by construction.
\item Cobb--Douglas overstates aggregate markups by 8.5--16 pp across countries. US manufacturing: +6.6 pp (1.35 in the 1960s to 1.41 in the 2010s) vs +18 pp under Cobb--Douglas.
\end{itemize}
\item \alert{Raval's flexible cost shares} (2019 WP; Raval 2022): bin producers by $w^L L/(w^M M)$, a signal of $\omega^L$, and use cost shares within bins. Labor and materials markups become positively correlated; markups rise 2--9\% per doubling of sales.
\item \alert{Doraszelski and Jaumandreu (2019)}: translog with both productivities, estimated jointly with the markup (dynamic panel or OP/LP). Exporters and non-exporters have about the same markups, and DLW-style markups are biased up.
\item Raval's analogy: demand estimation learned to allow heterogeneity (random coefficients). The production approach needs heterogeneity in technology.
\end{itemize}
\end{frame}


\begin{frame}{Every Wedge Looks Like a Markup}
The production approach identifies $\theta^X/\alpha^X$: the markup \alert{times} any other wedge in the FOC for $X$.
\begin{itemize}
\item Adjustment costs: $\theta^X/\alpha^X = \mu\,(1 + \tau^X)$ (DLW fn.\ 50); $\mu\,(1 + \kappa')$ (Bond et al.\ 2021).
\item Input market power (Rubens 2023), with markdown $\psi^V = 1 + $ inverse input-supply elasticity:
\begin{align*}
\mu = \Big(\psi^L \frac{\alpha^L}{\theta^L} + \psi^M \alpha^M\Big)^{-1},
\end{align*}
which nests $\mu = \theta^L/\alpha^L$ (DLW), $\mu = \theta^L/(\alpha^L \psi^L)$ (Morlacco), and $\mu = (\alpha^L/\theta^L + \alpha^M)^{-1}$ (Leontief, no markdowns).
\item Production data alone cannot split a markup from a markdown. Rubens adds an input-supply model; the alternative is a downstream demand model.
\item Rubens's application: a 2002 order closed small Chinese cigarette makers (340 $\to$ 167). Leaf's revenue share fell from 0.41 to 0.28 while labor's stayed near 0.03, consistent with rising oligopsony power over farmers.
\end{itemize}
\end{frame}


\begin{frame}{What Each Approach Assumes}
\scriptsize
\renewcommand{\arraystretch}{1.25}
\begin{tabularx}{\textwidth}{@{}l l X X@{}}
\toprule
Method & $\omega_{t-1}$ replaced by & Rules out & Allows \\
\midrule
Cost shares & --- & adjustment costs; $\mu \ne 1$ (revenue shares) or non-CRS (cost shares) & any $f$, any demand (cost-share version) \\
Dynamic panel & lagged production function & nonlinear $g$; makes selection hard & unobserved cost and demand shocks; any conduct \\
OP & investment demand & capital-cost heterogeneity; wage and demand shocks & selection; nonparametric $g$ \\
LP & materials demand & wage and demand shocks; first stage fails (functional dependence) & zero investment \\
ACF & conditional materials demand & unobserved $w^M$; demand shocks & dynamic labor; wage shocks \\
GNR & FOC (share regression) & $\mu \neq 1$ unless demand is modeled; wedges in $m$ & gross output; nonparametric $f$ \\
Doraszelski--Li & step 1 with $k_{t+1}$ & very nonlinear $g$ with large $\operatorname{Var}(\zeta)$ & demand shocks; unobserved rivals \\
Jaumandreu & cost-share FOC ($mc = AVC/\nu$) & nonlinear $g$; FOC wedges & any conduct \\
Demirer & ratio of input demands + control variable & non-homothetic $(L,M)$; demand shocks & labor-augmenting productivity \\
\bottomrule
\end{tabularx}
\end{frame}


\section{Aggregation and Misallocation}

\begin{frame}{The Rise of Market Power?}
\begin{itemize}
\item De Loecker, Eeckhout and Unger (2020): the sales-weighted average firm markup in Compustat is 1.34 in the 1960s, 1.21 in 1980 and 1.61 in 2016. The median is flat and the upper tail rises; about two-thirds of the rise is reallocation toward high-markup firms. DLW with COGS as the variable input, sector-year Cobb--Douglas.
\item Every piece of this lecture bears on that number:
\begin{itemize}
\item \alert{which costs}: COGS alone, or COGS plus selling, general and administrative costs (Traina 2018; next slide);
\item \alert{which input}: labor, materials, or COGS give opposite trends (Raval 2023);
\item \alert{which technology}: Cobb--Douglas vs heterogeneous technology, +18 vs +6.6 pp in US manufacturing (Demirer 2026);
\item \alert{which identification}: known returns to scale give cost-weighted Compustat markups of about 1.04--1.20 (Flynn--Gandhi--Traina 2019);
\item \alert{which data}: revenue data cannot pin down levels, but trends can survive (Bond et al.\ 2021; De Ridder et al.\ 2026);
\item \alert{which average}: Edmond--Midrigan--Xu (below).
\end{itemize}
\end{itemize}
\end{frame}


\begin{frame}{Which Costs Are Variable? (Traina 2018; DLEU 2020)}
\small
\begin{columns}[T]
\begin{column}{0.52\textwidth}
\includegraphics[width=\textwidth]{traina_fig2_cogs_opex.png}

{\scriptsize Sales-weighted Compustat markups with COGS (red) or operating expenses, COGS + SG\&A (blue), as the variable input. Traina (2018), Fig.\ 2.}
\end{column}
\begin{column}{0.46\textwidth}
\begin{itemize}
\item DLEU's method, with operating expenses (COGS + SG\&A) as the variable input.
\item SG\&A rose from about 12\% to 22\% of operating expenses. A constant omitted share biases the \alert{level}; a \alert{rising} one biases the \alert{trend}.
\item COGS: 1.15 (1980) to about 1.40. Operating expenses: about 1.15 (1950), under 1.10 (early 1980s), back to the 1950 level. Reweighting to the Census size distribution makes it flat or falling.
\item \alert{DLEU's reply}: on their sample his measure still rises, 1.08 to 1.18 (1980--2016); with SG\&A as an input the markup rises from about 1.0 to 1.3; the profit rate rises from 1\% to 8\% while overhead goes from 15\% to 21\% of total cost.
\end{itemize}
\end{column}
\end{columns}
\end{frame}


\begin{frame}{Which Average Markup? (Edmond, Midrigan and Xu 2023)}
The markup that distorts aggregate input choices is the \alert{cost-weighted} mean. With constant returns to the variable bundle, $\mu_{jt} = R_{jt}/VC_{jt}$; with sales shares $s_{jt} = R_{jt}/\sum_k R_{kt}$:
\begin{align*}
\mathcal{M}_t \equiv \sum_j \mu_{jt}\, \frac{VC_{jt}}{\sum_k VC_{kt}} = \frac{\sum_j R_{jt}}{\sum_j VC_{jt}} = \Bigg(\sum_j \frac{s_{jt}}{\mu_{jt}}\Bigg)^{-1} \le \sum_j s_{jt}\, \mu_{jt}.
\end{align*}
\begin{itemize}
\item $\mathcal{M}_t$ is aggregate revenue over aggregate variable cost: the wedge in the aggregate FOC. The steps use only $\theta^X = \mu\,\alpha^X$, so they hold for any demand system or conduct (their App.\ A).
\item The inequality is harmonic $\le$ arithmetic mean, with equality only if markups are equal. The sales-weighted mean overweights high-markup firms, so it overstates the aggregate distortion whenever markups are dispersed.
\item Compustat 2012: 1.26 cost-weighted (1.21 counting SG\&A as cost) vs about 1.6 sales-weighted.
\end{itemize}
\end{frame}


\begin{frame}{Cost-Weighted vs Sales-Weighted Markups, Compustat}
\begin{center}
\includegraphics[width=0.7\textwidth]{emx_fig8_weights.png}
\end{center}
{\small Edmond, Midrigan and Xu (2023), Figure 8 (NBER w24800 version). The sales-weighted average rises from about 1.17 in 1980 to about 1.55; the cost-weighted average is flat near 1.25 after 2000.}
\end{frame}


\begin{frame}{How Costly Are Markups? (Edmond, Midrigan and Xu 2023)}
Kimball demand: larger, more productive firms face less elastic demand and charge higher markups.
\begin{itemize}
\item Benchmark aggregate markup 1.15 (Barkai 2017). Removing all markup distortions is worth \alert{7.5\%} of consumption:
\begin{itemize}
\item about 2/3 from the aggregate markup, which acts like a uniform output tax;
\item about 1/3 from misallocation (a TFP loss of 1.2\%, against 8.4\% under CES with the same aggregate markup);
\item almost nothing from inefficient entry.
\end{itemize}
\item With the Compustat cost-weighted markup (1.25): 18.6\%, four-fifths from the aggregate markup.
\item Policy:
\begin{itemize}
\item entry subsidies barely move $\mathcal{M}$, because direct and compositional effects offset;
\item size-dependent taxes lower $\mathcal{M}$ but worsen misallocation: the distorted economy has \emph{too little} concentration.
\end{itemize}
\item DLW markups on Taiwanese producers discipline the key demand parameter (how fast elasticities fall with size).
\end{itemize}
\end{frame}


\begin{frame}{Misallocation: Hsieh and Klenow (2009)}
\begin{itemize}
\item Benchmark: CES demand, CRS Cobb--Douglas technology, common input prices, no adjustment costs. Marginal revenue products are equal across producers, and $\text{TFPR}_{jt} = \ln\mu + \ln c(\mathbf{w}_t) - \ln \bar p_t$ is the same for every firm (shown earlier).
\item So any TFPR dispersion is read as wedges. With an output wedge $\tau^Y_{jt}$ and a capital wedge $\tau^K_{jt}$:
\begin{align*}
\text{TFPR}_{jt} = \text{const}_t + \beta_k \ln\big(1 + \tau^K_{jt}\big) - \ln\big(1 - \tau^Y_{jt}\big).
\end{align*}
\item Equalizing marginal products to the US level would raise manufacturing TFP by 30--50\% in China and 40--60\% in India.
\item But outside the benchmark TFPR also varies with demand and markups (Haltiwanger, Kulick and Syverson 2018; FHS 2008), and with the timing of inputs (next slide).
\end{itemize}
\end{frame}


\begin{frame}{\ldots or Dynamics? (Asker, Collard-Wexler and De Loecker 2014)}
Capital is chosen a period ahead. After optimizing out the flexible inputs, log revenue is $r_{jt} = c_t + \gamma\, k_{jt} + \tilde\omega_{jt}$, with $\gamma < 1$ and revenue productivity $\tilde\omega_{jt} = \rho\, \tilde\omega_{jt-1} + \xi_{jt}$. With no distortions the firm sets $\E_{t-1}[\text{MRPK}_{jt}] = w^K_t$, so
\begin{align*}
\ln \text{MRPK}_{jt} = \ln\gamma + c_t - (1 - \gamma)\, k_{jt} + \tilde\omega_{jt}
\quad\Longrightarrow\quad
\ln \text{MRPK}_{jt} = \ln w^K_t - \ln \E\big[e^{\xi}\big] + \xi_{jt}.
\end{align*}
\begin{itemize}
\item MRPK is dispersed with \alert{no} distortion: its cross-sectional variance is $\sigma^2_\xi$, the volatility of revenue productivity, which ACDL measure from a revenue production function. Adjustment costs add more.
\item Measured volatility and MRPK dispersion line up across industries in 8 countries and across 33 countries in the World Bank Enterprise Surveys.
\item A dynamic investment model with adjustment costs fits the cross-industry and cross-country dispersion with an uncentered $R^2$ of about 0.8--0.9.
\item Lesson: static marginal-product gaps are weak evidence of policy distortions. The production function, the timing of inputs, and the demand system all matter.
\end{itemize}
\end{frame}


\begin{frame}{Takeaways}
\begin{itemize}
\item Be explicit about \alert{timing}, the \alert{scalar unobservable}, and what moves a flexible input besides $\omega$. That last item is your exclusion restriction.
\item Use ACF-style conditional proxies rather than LP's first stage. For gross output use GNR (price takers) or a Leontief argument, and cross-check against dynamic panel.
\item Under market power, put demand and conduct shifters (and $k_{t+1}$) into step 1, or use cost-share forms of the FOC.
\item Put firm actions (exporting, R\&D) into the law of motion $g$, not into a second-stage regression on $\hat\omega$.
\item Markups:
\begin{itemize}
\item levels need output prices; dispersion and trends survive revenue data if technology is common;
\item check that different inputs agree, and allow labor-augmenting technology when they do not;
\item every wedge (adjustment costs, markdowns) looks like a markup;
\item aggregate with cost weights (harmonic sales weights) for welfare.
\end{itemize}
\item The production approach complements the demand approach. It measures market power across many industries; it cannot simulate a merger.
\end{itemize}
\end{frame}

\end{document}
